\appendix{Представление графического материала}

Графический материал, выполненный на отдельных листах,
изображен на рисунках А.1--А.\arabic{числоПлакатов}.
\setcounter{числоПлакатов}{0}

\renewcommand{\thefigure}{А.\arabic{figure}} % шаблон номера для плакатов

\begin{landscape}
	
	\begin{плакат}
		
\includegraphics[width=0.82\linewidth]{Svedenia.eps}
		\заголовок{Сведения о ВКРБ}
		\label{pl1:image}
	\end{плакат}
	
	\begin{плакат}
		\includegraphics[width=0.82\linewidth]{Tseli.eps}
		\заголовок{Цель и задачи разработки}
		\label{pl2:image}
	\end{плакат}
	
	\begin{плакат}
		\includegraphics[width=0.82\linewidth]{plakat-component-diagram.eps}
		\заголовок{Диаграмма компонентов приложения}
		\label{pl3:image}
	\end{плакат}
	
	\begin{плакат}
		
\includegraphics[width=0.82\linewidth]{plakat-data-flow-app.eps}
		\заголовок{Диаграмма потока данных в приложении}
		\label{pl4:image}
	\end{плакат}
	
	\begin{плакат}
		
\includegraphics[width=0.82\linewidth]{plakat-er.eps}
		\заголовок{Схема базы данных программной системы}
		\label{pl5:image}
	\end{плакат}
	
	\begin{плакат}
		
\includegraphics[width=0.82\linewidth]{plakat-workflow.eps}
		\заголовок{Жизненный цикл заявки}
		\label{pl6:image}
	\end{плакат}
	
	\begin{плакат}
		\includegraphics[width=0.82\linewidth]{plakat-interaction.eps}
		\заголовок{Диаграмма взаимодействия компонентов}
		\label{pl7:image}
	\end{плакат}
	
	\begin{плакат}
		\includegraphics[width=0.82\linewidth]{plakat-project-structure.eps}
		\заголовок{Структура программного проекта}
		\label{pl8:image}
	\end{плакат}
	
	\begin{плакат}
		\includegraphics[width=0.82\linewidth]{Zaklyuchenie.eps}
		\заголовок{Заключение}
		\label{pl9:image}
	\end{плакат}
	
\end{landscape}